% =====================================================================================
% Section 6 -- Conclusion
% =====================================================================================
\section{Conclusion}
\label{sec:conclusion}

Do System One models make better rerankers? On the same backbone, no: reranker fine-tuning beats decision fine-tuning
on every backbone where both exist, also on collections without training data, and decision fine-tuning barely
improves on the untuned base model (RQ1). Nor are decision models at least better calibrated: reranker fine-tunes
of the same backbone are about as well calibrated, give the better set at a fixed cutoff, and a reranker fine-tune
with a pointwise loss on the yes/no logit is calibrated better than any decision model (RQ2). For the IR community,
the probabilities that set decision models apart are a property of the training loss, available to any reranker
fine-tune.

\paragraph{Limitations.}
No reranker fine-tune of the 27B backbone exists yet, so at 27B the comparison is open. The fine-tunes saw different
training data (MS~MARCO for the reranker fine-tunes, not for Kev), which the zero-shot subset only partly controls;
benchmark contamination of the backbones and of \jev{} cannot be ruled out. Literature numbers are not recomputed
and not tested; we report no significance tests. Calibration is measured on TREC DL only, the only collections with
judged non-relevant passages. Other uses of calibrated probabilities, such as query performance prediction, are left
to future work.
